\documentclass[aspectratio=169]{beamer}

\usetheme{Madrid}
\definecolor{unsblue}{RGB}{30,60,120}
\definecolor{unslightblue}{RGB}{220,230,245}
\definecolor{alertred}{RGB}{180,30,30}
\usecolortheme[named=unsblue]{structure}
\setbeamercolor{palette primary}{bg=unsblue, fg=white}
\setbeamercolor{palette secondary}{bg=unsblue!80, fg=white}
\setbeamercolor{palette tertiary}{bg=unsblue!60, fg=white}
\setbeamercolor{palette quaternary}{bg=unsblue, fg=white}
\setbeamercolor{block title}{bg=unsblue, fg=white}
\setbeamercolor{block body}{bg=unslightblue}
\setbeamercolor{block title alerted}{bg=alertred, fg=white}
\setbeamercolor{block body alerted}{bg=red!8}
\setbeamercolor{block title example}{bg=unsblue!60!black, fg=white}
\setbeamercolor{block body example}{bg=unsblue!8}

\usepackage[utf8]{inputenc}
\usepackage[T1]{fontenc}
\usepackage{babel}
\babelprovide[main, import]{spanish}
\usepackage{amsmath, amssymb}
\usepackage{bm}
\usepackage{booktabs}

\newcommand{\R}{\mathbb{R}}
\newcommand{\bx}{\bm{x}}
\newcommand{\bw}{\bm{w}}

\title[MAD1 -- Unidad 4]{Aprendizaje Supervisado}
\subtitle{Clase 22: SVM y K-Nearest Neighbors}
\author[J. Bavio]{Dr.\ José Bavio}
\institute[UNS]{Departamento de Matemática\\Universidad Nacional del Sur}
\date{2026}

\begin{document}

\begin{frame}
  \titlepage
\end{frame}

\begin{frame}{Contenido de la clase}
  \tableofcontents
\end{frame}

% ═════════════════════════════════════════════════════════════════════════════
\section{Support Vector Machines}
% ═════════════════════════════════════════════════════════════════════════════

\begin{frame}{Idea central de SVM}
  Problema: clasificar datos en dos clases ($y \in \{-1, +1\}$).\\
  Hay muchos hiperplanos que separan las clases. ¿Cuál elegir?

  \medskip
  \begin{block}{Respuesta de SVM: el de máximo margen}
    Elegir el hiperplano que maximiza la \textbf{distancia} entre sí mismo
    y las observaciones más cercanas de cada clase.
  \end{block}

  \medskip
  \begin{exampleblock}{Intuición geométrica}
    Es como trazar la línea divisoria entre dos territorios
    de forma que quede lo más lejos posible de ambos lados.\\
    Cuanto más ``espacio libre'' hay a ambos lados, más confianza tenemos
    al clasificar un punto nuevo que cae cerca del borde.
  \end{exampleblock}
\end{frame}

\begin{frame}{Hiperplano y margen}
  El hiperplano se define como:
  \[
    \mathcal{H} = \{\bx : \bw^\top \bx + b = 0\}
  \]

  Los \textbf{vectores de soporte} son las observaciones más cercanas al hiperplano:
  las únicas que determinan la solución. El resto de los datos son irrelevantes.

  \medskip
  \begin{exampleblock}{Coloquialmente}
    De todos los datos de entrenamiento, solo importan los que están ``al borde''.\\
    Si movés cualquier otro punto, el hiperplano no cambia.\\
    Si movés un vector de soporte, sí cambia.
  \end{exampleblock}

  \medskip
  Margen $= 2/\|\bw\|$. Maximizar el margen $\Leftrightarrow$ minimizar $\|\bw\|^2$.
\end{frame}

\begin{frame}{SVM de margen duro}
  \begin{block}{Problema de optimización (datos separables)}
    \[
      \min_{\bw, b} \;\frac{1}{2}\|\bw\|^2
      \quad\text{sujeto a}\quad
      y_i(\bw^\top \bx_i + b) \geq 1 \quad \forall\, i
    \]
  \end{block}

  \medskip
  La restricción $y_i(\bw^\top \bx_i + b) \geq 1$ dice:
  cada observación debe estar al lado correcto del margen (no solo del hiperplano).

  \medskip
  \begin{alertblock}{Limitación: datos no separables}
    Si las clases se solapan, ningún hiperplano satisface todas las restricciones.
    El margen duro no existe. Hay que relajar las restricciones.
  \end{alertblock}
\end{frame}

\begin{frame}{SVM de margen blando}
  Se introducen \textbf{variables de holgura} $\xi_i \geq 0$ que permiten violaciones:
  \[
    \min_{\bw, b, \bm{\xi}} \;\frac{1}{2}\|\bw\|^2 + C\sum_{i=1}^{n} \xi_i
    \quad\text{sujeto a}\quad
    y_i(\bw^\top \bx_i + b) \geq 1 - \xi_i,\;\; \xi_i \geq 0
  \]

  \medskip
  El parámetro $C$ controla el compromiso:
  \begin{itemize}
    \item $C$ \textbf{grande}: penaliza mucho las violaciones $\Rightarrow$ margen más angosto, más riesgo de sobreajuste.
    \item $C$ \textbf{pequeño}: tolera violaciones $\Rightarrow$ margen más amplio, más regularización.
  \end{itemize}

  \medskip
  \begin{exampleblock}{Analogía}
    $C$ es como el nivel de tolerancia del modelo.
    Un $C$ alto dice: ``no acepto ningún error''. Un $C$ bajo dice: ``acepto algunos errores si eso me da más margen general''.
  \end{exampleblock}
\end{frame}

\begin{frame}{El truco del kernel}
  ¿Qué pasa si los datos no son linealmente separables en ninguna dimensión razonable?

  \medskip
  \begin{block}{Idea: mapear a mayor dimensión}
    Proyectar $\bx$ a un espacio de alta dimensión $\phi(\bx)$ donde las clases sí sean separables.\\
    El hiperplano en ese espacio corresponde a una frontera \textbf{no lineal} en el espacio original.
  \end{block}

  \medskip
  \begin{exampleblock}{Truco del kernel}
    La formulación dual de SVM solo usa $\bx_i$ a través de productos internos $\langle \bx_i, \bx_j \rangle$.\\
    Si reemplazamos $\langle \bx_i, \bx_j \rangle$ por $K(\bx_i, \bx_j) = \langle \phi(\bx_i), \phi(\bx_j) \rangle$,
    calculamos el producto interno en el espacio transformado \textbf{sin calcularlo explícitamente}.\\[4pt]
    Es como operar en una dimensión altísima sin pagar el costo computacional.
  \end{exampleblock}
\end{frame}

\begin{frame}{Kernels más usados}
  \begin{center}
  \begin{tabular}{lll}
    \toprule
    \textbf{Kernel} & \textbf{Fórmula} & \textbf{Cuándo usarlo} \\
    \midrule
    Lineal     & $\bx^\top \bx'$                         & Datos de alta dim.\ (texto) \\
    Polinomial & $(\gamma\,\bx^\top\bx' + r)^d$          & Interacciones polinomiales \\
    RBF        & $\exp(-\gamma\|\bx-\bx'\|^2)$           & Uso general, no lineal \\
    Sigmoide   & $\tanh(\gamma\,\bx^\top\bx' + r)$       & Similar a redes neuronales \\
    \bottomrule
  \end{tabular}
  \end{center}

  \medskip
  \begin{block}{Kernel RBF: el parámetro $\gamma$}
    \begin{itemize}
      \item $\gamma$ pequeño: fronteras suaves, más regularización.
      \item $\gamma$ grande: fronteras muy ajustadas al entrenamiento, riesgo de sobreajuste.
    \end{itemize}
  \end{block}
\end{frame}

\begin{frame}{SVM multiclase y consideraciones prácticas}
  \begin{block}{SVM es binaria por definición}
    Para $K > 2$ clases se usan estrategias de descomposición:
    \begin{itemize}
      \item \textbf{Uno vs.\ todos (OvR):} un clasificador por clase.
      \item \textbf{Uno vs.\ uno (OvO):} un clasificador por par de clases; predicción por votación.
    \end{itemize}
  \end{block}

  \medskip
  \begin{alertblock}{Escala de los datos}
    SVM es sensible a la escala de las variables. Siempre estandarizar antes de entrenar.
    No hacerlo es uno de los errores más comunes en la práctica.
  \end{alertblock}

  \medskip
  \begin{exampleblock}{¿Cuándo usar SVM?}
    Funciona muy bien con pocos datos y alta dimensión (texto, imágenes pequeñas).
    Costoso para $n$ muy grande (la solución requiere todos los vectores de soporte).
  \end{exampleblock}
\end{frame}

% ═════════════════════════════════════════════════════════════════════════════
\section{K-Nearest Neighbors}
% ═════════════════════════════════════════════════════════════════════════════

\begin{frame}{La idea detrás de KNN}
  \begin{block}{Principio}
    Para clasificar un punto nuevo $\bx_0$:
    encontrar sus $k$ vecinos más cercanos en el conjunto de entrenamiento
    y asignar la clase más frecuente entre ellos.
  \end{block}

  \medskip
  \begin{exampleblock}{Analogía cotidiana}
    ``Dime con quién andas y te diré quién eres.''\\
    KNN asume que los puntos cercanos en el espacio de features pertenecen
    a la misma clase. Si los $k$ vecinos más cercanos son 4 de clase A y 1 de clase B,
    predecimos clase A.
  \end{exampleblock}

  \medskip
  \textbf{No hay entrenamiento explícito:} el modelo ``aprende'' guardando todos los datos.
  La predicción es el paso costoso.
\end{frame}

\begin{frame}{Definición formal}
  \begin{block}{Clasificador $k$-NN}
    \[
      \hat{y}(\bx_0) = \arg\max_c \sum_{i \in \mathcal{N}_k(\bx_0)} \mathbf{1}(y_i = c)
    \]
    donde $\mathcal{N}_k(\bx_0)$ es el conjunto de los $k$ índices con menor distancia a $\bx_0$.
  \end{block}

  \medskip
  La distancia habitual es la \textbf{euclídea}:
  \[
    d(\bx, \bx') = \sqrt{\sum_{j=1}^{p}(x_j - x'_j)^2}
  \]
  Alternativas: Manhattan ($L_1$), Minkowski ($L_p$), Mahalanobis.
\end{frame}

\begin{frame}{Elección de $k$: el hiperparámetro clave}
  \begin{columns}
    \column{0.48\textwidth}
    \begin{alertblock}{$k$ pequeño ($k=1$)}
      Frontera muy irregular.\\
      El modelo memoriza cada dato.\\
      Alta varianza, bajo sesgo.\\[4pt]
      \textit{Como adivinar el partido de toda la escuela a partir del voto de una sola persona.}
    \end{alertblock}

    \column{0.48\textwidth}
    \begin{alertblock}{$k$ grande}
      Frontera muy suave.\\
      El modelo ignora detalles locales.\\
      Bajo sesgo, alto sesgo global.\\[4pt]
      \textit{Como preguntarle a toda la Argentina para clasificar un barrio.}
    \end{alertblock}
  \end{columns}

  \medskip
  \begin{block}{En la práctica}
    Se elige $k$ por validación cruzada.
    Probar valores impares para evitar empates en clasificación binaria.
    Valores típicos: $k = 5, 7, 11$.
  \end{block}
\end{frame}

\begin{frame}{La maldición de la dimensionalidad}
  \begin{block}{El problema}
    En espacios de alta dimensión ($p$ grande), las distancias entre puntos
    tienden a homogeneizarse: todos los puntos quedan ``igual de lejos''.
  \end{block}

  \medskip
  \begin{exampleblock}{Un dato sorprendente}
    En $p = 1$ dimensión, el vecino más cercano está a distancia $\sim 1/n$ en promedio.\\
    En $p = 100$ dimensiones, para tener el 1\% de los datos ``cerca'',
    necesitaría cubrir el 63\% del rango de cada variable.\\[4pt]
    El espacio se vuelve enormemente vacío conforme $p$ crece.
  \end{exampleblock}

  \medskip
  \begin{alertblock}{Consecuencia para KNN}
    El concepto de ``vecino cercano'' pierde significado en alta dimensión.
    KNN funciona bien con $p$ pequeño (idealmente $\leq 20$).
  \end{alertblock}
\end{frame}

\begin{frame}{Consideraciones prácticas de KNN}
  \begin{alertblock}{Escala: obligatorio normalizar}
    Si $x_1 \in [0, 1]$ y $x_2 \in [0, 10{.}000]$,
    la distancia euclídea queda dominada por $x_2$.
    $x_1$ queda prácticamente invisible para el modelo.\\
    \textbf{Siempre estandarizar o normalizar antes de KNN.}
  \end{alertblock}

  \medskip
  \begin{block}{Costo computacional}
    \begin{itemize}
      \item Entrenamiento: prácticamente nulo (solo guardar los datos).
      \item Predicción: $\mathcal{O}(n \cdot p)$ por cada punto a predecir.
      \item Con $n$ grande, la predicción se vuelve muy lenta.
    \end{itemize}
  \end{block}

  \medskip
  \begin{exampleblock}{¿Cuándo usar KNN?}
    Pocos datos, dimensión baja, y cuando el modelo necesita ser simple
    y no importa el tiempo de predicción. También es útil como baseline.
  \end{exampleblock}
\end{frame}

% ─────────────────────────────────────────────────────────────────────────────
\section{Comparación SVM vs KNN}
% ─────────────────────────────────────────────────────────────────────────────

\begin{frame}{SVM vs KNN: resumen comparativo}
  \begin{center}
  \begin{tabular}{lll}
    \toprule
    \textbf{Aspecto} & \textbf{SVM} & \textbf{KNN} \\
    \midrule
    Tipo de frontera & Global (hiperplano + kernel) & Local (mayoría de vecinos) \\
    Entrenamiento    & Costoso ($n$ grande)         & Instantáneo \\
    Predicción       & Rápida                        & Lenta ($n$ grande) \\
    Escala           & Obligatorio normalizar        & Obligatorio normalizar \\
    Alta dimensión   & Bien (kernel lineal en texto) & Mal (maldición dim.) \\
    Interpretación   & Vectores de soporte           & Directa pero local \\
    Hiperparámetros  & $C$, kernel, $\gamma$         & $k$, métrica \\
    \bottomrule
  \end{tabular}
  \end{center}
\end{frame}

\begin{frame}{Resumen de la clase}
  \begin{itemize}
    \item \textbf{SVM} busca el hiperplano de máximo margen entre clases.
    \item El \textbf{margen blando} (parámetro $C$) permite violaciones controladas.
    \item El \textbf{truco del kernel} extiende SVM a fronteras no lineales sin calcular explícitamente la transformación.
    \item \textbf{KNN} predice por mayoría entre los $k$ vecinos más cercanos.
    \item Elegir $k$ es el compromiso sesgo-varianza del método.
    \item La \textbf{maldición de la dimensionalidad} limita KNN en espacios de alta dimensión.
    \item \textbf{Ambos modelos requieren normalización previa.}
  \end{itemize}

  \bigskip
  \textbf{Próxima clase:} Redes neuronales, métricas de evaluación, SMOTE y pipelines.
\end{frame}

\end{document}
